% GENERATED FILE. Edit canonical lessons, then run npm run generate:books.
% canonical-source-hash: fnv1a64:c30859682d5d2fed
% canonical-lessons: GE-C29-freund, GE-C29-freundin, GE-C29-familie

\chapter{Friend and Family}
\label{ch:ge-friend-and-family}
\hlchaptermodality{29}

\begin{chapteropening}
\textbf{By the end of this chapter:} \emph{I can name a friend, male or female, and my family in German, build a feminine person-noun with -in, and say which of the three words is native and which is borrowed.}

The chapter closes on Familie, the one Latin loan beside two native words: the reader produces der Freund, die Freundin and die Familie with articles, builds the feminine -in suffix on a masculine person-noun and names its one English fossil (vixen), and states that Familie gathers the very members Chapter 10 already named (Vater, Mutter, Bruder, Schwester, Geschwister). It reaches back to Chapter 2's freut mich, which shares Freund's root, and to Chapter 10's family words.
\end{chapteropening}

\section[der Freund]{der Freund — a word that was already a compliment}

\label{lesson:GE-C29-freund}



\begin{quote}
Chapter 2 taught you \textbf{freut mich} — \textquotedblleft{}pleased to meet you,\textquotedblright{} literally \textquotedblleft{}it gladdens me.\textquotedblright{} The person doing the gladdening has a name of their own, and it is built from the very same root.
\end{quote}

\subsection*{You'll want to know: der Freund}
\begin{quote}
\textbf{der Freund} — friend (male, or a friend in general). Masculine.
\end{quote}

\begin{quote}
\textbf{Das ist der Freund.} — \textquotedblleft{}This/That is the friend.\textquotedblright{}
\end{quote}

\begin{sounds}
\textbf{eu} is a diphthong said \emph{oy}, exactly as in English \emph{boy}: \textbf{Freund} = \emph{FROYND}. It is the same sound spelled \textbf{äu} in \emph{Häuser} and \emph{läuft}, just without the umlaut in front of it.
\end{sounds}

\begin{cousinweb}
\textbf{Freund} and English \textbf{friend} are not merely cousins — they are the exact same inherited word, Proto-Germanic \emph{\textbf{*frijōnds}}. That word was never a plain noun to begin with: it is a \textbf{frozen present participle} of the verb \emph{\textbf{*frijōną}}, \textquotedblleft{}to love,\textquotedblright{} so \emph{Freund} literally means \textquotedblleft{}the one who is loving.\textquotedblright{} \emph{freut mich} is built on the same family: both trace to Proto-Indo-European \emph{\textbf{*preyH-}}, \textquotedblleft{}to love, to please.\textquotedblright{} The goddess \textbf{Frigg} — whose day German keeps as \emph{Freitag} and English as \textbf{Friday} — carries the same root one step further.

\emph{Milch} needed no journey to reach German; \emph{Freund} did not either — but where \emph{Milch} is a noun that was always a noun, \emph{Freund} is a verb that stopped moving and became one.
\end{cousinweb}

\subsection*{Your turn}
Say these aloud:

\begin{itemize}
\raggedright
  \item \textquotedblleft{}der Freund\textquotedblright{} — \emph{FROYND}
  \item \textquotedblleft{}Das ist der Freund.\textquotedblright{}
  \item the family — \textquotedblleft{}Freund, freut mich, Frigg, Friday — all one root\textquotedblright{}
\end{itemize}

\begin{tcolorbox}[breakable,colback=teal!4,colframe=teal!35!black,title={Before you move on}]
Give \textquotedblleft{}friend\textquotedblright{} with its article. (\textbf{Der Freund}.) What part of speech did \emph{Freund} start out as? (\textbf{A verb} — a frozen present participle of \textquotedblleft{}to love.\textquotedblright{}) Which Chapter 2 phrase shares its root? (\textbf{Freut mich}.) Which day of the week carries the same PIE root one step further? (\textbf{Freitag / Friday}, from the goddess \textbf{Frigg}.)
\end{tcolorbox}
\section[die Freundin]{die Freundin — one suffix, and the fox that kept it in English}

\label{lesson:GE-C29-freundin}



\begin{quote}
You have \textbf{der Freund}. German does not need a new word for the feminine — it needs one small, reusable ending.
\end{quote}

\subsection*{You'll want to know: die Freundin}
\begin{quote}
\textbf{die Freundin} — friend (female). Feminine.
\end{quote}

\begin{quote}
\textbf{Das ist die Freundin.} — \textquotedblleft{}This/That is the (female) friend.\textquotedblright{}
\end{quote}

\begin{grammarlens}[title={Grammar Lens: the ending that builds a whole class of nouns}]
Add \textbf{-in} to a masculine noun for a person and you get the feminine: \textbf{der Lehrer $\to$ die Lehrerin} (\textquotedblleft{}teacher\textquotedblright{}), \textbf{der Student $\to$ die Studentin} (\textquotedblleft{}student\textquotedblright{}), \textbf{der Freund $\to$ die Freundin}. It is entirely regular, and it works on almost any such noun you will ever meet — one rule, not a new word each time.
\end{grammarlens}

\begin{cousinweb}
\textbf{-in} is native Germanic, not a Latin import — English inherited the same suffix and mostly stopped using it. Almost mostly: one English word still carries it, fossilized, unrecognized as a suffix by most speakers who use it — \textbf{vixen}, \textquotedblleft{}female fox,\textquotedblright{} built the same way \emph{Freundin} is built. English did not lose this pattern by borrowing over it; it simply stopped reaching for it, and later borrowed French \emph{-ess} (\textbf{princess, actress}) to do the same job instead. German never stopped.
\end{cousinweb}

\subsection*{Your turn}
Say these aloud:

\begin{itemize}
\raggedright
  \item \textquotedblleft{}die Freundin\textquotedblright{} — Freund plus -in
  \item \textquotedblleft{}Das ist die Freundin.\textquotedblright{}
  \item the pattern — \textquotedblleft{}Lehrer/Lehrerin, Student/Studentin, Freund/Freundin\textquotedblright{}
  \item the English fossil — \textquotedblleft{}fox, vixen — the same -in, once\textquotedblright{}
\end{itemize}

\begin{tcolorbox}[breakable,colback=teal!4,colframe=teal!35!black,title={Before you move on}]
Give \textquotedblleft{}friend (female)\textquotedblright{} with its article. (\textbf{Die Freundin}.) What suffix builds it from \emph{Freund}, and what does that suffix do generally? (\textbf{-in} — it makes the feminine of a person-noun.) Is \emph{-in} borrowed or native? (\textbf{Native} — English inherited it too.) Name the one English word that still shows it. (\textbf{Vixen.})
\end{tcolorbox}
\section[die Familie]{die Familie — the word that gathers Eltern and Geschwister}

\label{lesson:GE-C29-familie}



\begin{quote}
You have \textbf{der Freund} and \textbf{die Freundin}, both native, both built on a verb meaning \textquotedblleft{}to love.\textquotedblright{} The word for the group at home has a completely different history.
\end{quote}

\subsection*{You'll want to know: die Familie}
\begin{quote}
\textbf{die Familie} — family. Feminine.
\end{quote}

\begin{quote}
\textbf{Das ist die Familie.} — \textquotedblleft{}This/That is the family.\textquotedblright{}
\end{quote}

Chapter 10 already gave you the members: \textbf{Vater, Mutter, Bruder, Schwester} — and the collective \textbf{Geschwister}, \textquotedblleft{}siblings.\textquotedblright{} \emph{Familie} is the word that gathers all of them into one group.

\begin{sounds}
\textbf{Familie} keeps all three vowels distinct: \emph{fa-MEE-lee-eh}, four syllables in careful speech, with the stress on the second. It does not collapse the way English \emph{family} does.
\end{sounds}

\begin{cousinweb}
\textbf{Familie} is a Latin loan: \textbf{familia}, \textquotedblleft{}household,\textquotedblright{} related to \textbf{famulus}, \textquotedblleft{}servant\textquotedblright{} or \textquotedblleft{}slave\textquotedblright{} — a Roman household word, not originally a word about blood relation at all. English \textbf{family} borrowed the very same Latin root, by way of Old French, well before German did; German took \emph{Familie} directly from Latin later, in the early modern period, which is why the two loans look so close to identical without one having been copied from the other.

So this small chapter runs the same shape twice: \textbf{Freund} and \textbf{Freundin} are native, built from a Germanic verb meaning \textquotedblleft{}to love\textquotedblright{}; \textbf{Familie} is borrowed, from a Roman word that first meant \textquotedblleft{}household staff.\textquotedblright{} Chapter 11 set up exactly this pattern with \textbf{Wasser} native beside \textbf{Wein} borrowed, and Chapter 28 just ran it a third time with \textbf{Milch} beside \textbf{Kaffee} and \textbf{Tee}.
\end{cousinweb}

\subsection*{Your turn}
Say these aloud:

\begin{itemize}
\raggedright
  \item \textquotedblleft{}die Familie\textquotedblright{} — \emph{fa-MEE-lee-eh}
  \item \textquotedblleft{}Das ist die Familie.\textquotedblright{}
  \item the members inside it — \textquotedblleft{}Vater, Mutter, Bruder, Schwester, Geschwister\textquotedblright{}
  \item native beside borrowed — \textquotedblleft{}Freund, Freundin native; Familie borrowed\textquotedblright{}
\end{itemize}

\begin{tcolorbox}[breakable,colback=teal!4,colframe=teal!35!black,title={Before you move on}]
Give \textquotedblleft{}family\textquotedblright{} with its article. (\textbf{Die Familie}.) Is it native or borrowed? (\textbf{Borrowed} — Latin \emph{familia}, \textquotedblleft{}household.\textquotedblright{}) What did \emph{familia} first mean, before \textquotedblleft{}relatives\textquotedblright{}? (\textbf{Household} — kin to \emph{famulus}, \textquotedblleft{}servant.\textquotedblright{}) Which two words from this chapter are native instead? (\textbf{Freund} and \textbf{Freundin}.)
\end{tcolorbox}
